\section{Introduction}

\label{sec:introduction}

The metastable consensus family achieves agreement through repeated stochastic sampling,
where each validator queries a small random subset of peers and adopts the majority
preference. Slush~\cite{rocket2019wave}, the simplest member, has no memory
and can be driven back and forth by an adversary. The Wave protocol~\cite{kelling2020wave}
introduces cumulative confidence counters that make reversal exponentially unlikely
but does not provide a clear \emph{finality} signal.

Flare bridges this gap by introducing \emph{conviction counters}: a counter $cnt$ that
tracks consecutive rounds in which the sampled supermajority agrees with the validator's
current preference. When $cnt$ reaches threshold $\beta_2$, the validator issues a
\emph{certificate}---a cryptographic attestation of finality. If the preference changes
before $\beta_2$ is reached, the counter resets to zero (a \emph{skip}).

The certificate/skip mechanism creates a sharp binary: either a validator builds sufficient
consecutive evidence to certify, or it starts over. This paper proves that the
``starting over'' event becomes exponentially unlikely once the honest population
has converged, and that the certificate/skip structure prevents conflicting certificates
from forming.

\paragraph{Contributions.}
\begin{enumerate}[nosep]
\item The Flare protocol specification with certificate/skip conviction counters and
DAG integration (\S\ref{sec:protocol}).
\item The Certificate Skip Exclusivity (CSE) theorem: conflicting certificates cannot
coexist under honest supermajority (\S\ref{sec:cse}).
\item Analysis of luring attacks and the honest support invariant (\S\ref{sec:attacks}).
\item A comparative analysis with Wave showing that Flare provides sharper finality
at the cost of slightly higher latency (\S\ref{sec:comparison}).
\item Empirical results from a 500-node testnet (\S\ref{sec:evaluation}).
\end{enumerate}

%% -----------------------------------------------------------------------
%%  2. SYSTEM MODEL
%% -----------------------------------------------------------------------
